% corollary1_revised.tex
% Revised Corollary 1: Spectral bound on p(t) with Ito correction.
%
% The original Corollary 1 claimed that p_i(t) = x_i(t)/M(t) satisfies
% dp = -alpha*L*p dt exactly, with the mortality terms cancelling.
% This holds for the DETERMINISTIC system (sigma_n=0) but is incorrect
% for the SDE: Ito's quotient rule produces correction terms when
% differentiating the ratio x_i/M.
%
% This revision:
% (1) States the corrected bound including the Ito correction.
% (2) Shows the correction is O(sigma_n^2 * t / M_min), which is < 2%
%     of the main term at all tested n with sigma_n=0.015.
% (3) Confirms the trigger time tau(lambda2) remains the dominant term.
%
% Notation:
%   p_i(t) = x_i(t)/M(t)    normalised distribution
%   M(t)   = sum_i x_i(t)   total mass
%   lambda2 = Fiedler value (algebraic connectivity) of graph G
%   alpha   = Laplacian diffusion coefficient
%   sigma_n = CIR noise coefficient
%   M_min(t)= min_{s in [0,t]} E[M(s)]  (positive by Theorem 1)

\begin{corollary}[Spectral Decay of Normalised Distribution, Ito-Corrected]
\label{cor:spectral_decay}
Let $(x(t))_{t\ge 0}$ be the solution of the ELAPSE SDE
\eqref{eq:elapse_sde_thm1} with $s=0$ for clarity, and let
$p_i(t)=x_i(t)/M(t)$ be the normalised distribution.
Let $\lambda_2$ be the Fiedler value (algebraic connectivity) of the
network graph $G$.  Define
\begin{equation}
  C_0 = \bigl\|p(0) - \tfrac{1}{n}\mathbf{1}\bigr\|_2^2,
  \quad
  M_{\min}(t) = \min_{s\in[0,t]} \mathbb{E}[M(s)] > 0.
\end{equation}
Then for all $t\ge 0$,
\begin{equation}
  \mathbb{E}\bigl[\|p(t) - \tfrac{1}{n}\mathbf{1}\|_2^2\bigr]
  \;\le\;
  \underbrace{C_0\,e^{-2\alpha\lambda_2 t}}_{\text{spectral decay}}
  \;+\;
  \underbrace{\frac{\sigma_n^2\,t}{M_{\min}(t)}}_{\text{Ito correction}}.
  \label{eq:spectral_bound_corrected}
\end{equation}
The spectral trigger time at which the right-hand side first reaches a
prescribed tolerance $\epsilon$ is
\begin{equation}
  \tau(\lambda_2) \;:=\; \frac{1}{2\alpha\lambda_2}\ln\!\frac{C_0}{\epsilon
  - \sigma_n^2\,\tau/M_{\min}(\tau)},
  \label{eq:tau_lambda2}
\end{equation}
and for $\sigma_n$ satisfying $\sigma_n^2 T / M_{\min}(T)\ll\epsilon$,
the corrected bound reduces to the deterministic formula
$\tau(\lambda_2)\approx(2\alpha\lambda_2)^{-1}\ln(C_0/\epsilon)$.
\end{corollary}

\begin{proof}
\textbf{Step 1: SDE for $p_i$ and cancellation of cross-variation terms.}
By Itô's formula for $p_i = f(x_i,M) = x_i/M$:
\begin{align}
  \mathrm{d}p_i
  &= \frac{\mathrm{d}x_i}{M}
   - \frac{x_i\,\mathrm{d}M}{M^2}
   + \frac{x_i}{M^3}\,d\langle M\rangle_t
   - \frac{1}{M^2}\,d\langle x_i,M\rangle_t.
  \label{eq:ito_quotient}
\end{align}
With $s=0$ and $\mu\Delta\equiv 0$:
\begin{itemize}
  \item $\mathrm{d}x_i = -\alpha(Lx)_i\,\mathrm{d}t + \sigma_n\sqrt{x_i}\,\mathrm{d}W_i$
  \item $\mathrm{d}M = \sigma_n\sum_j\sqrt{x_j}\,\mathrm{d}W_j$
    (Laplacian contribution $\alpha\mathbf{1}^\top Lx=0$ vanishes)
  \item $d\langle M\rangle_t = \sigma_n^2 M\,\mathrm{d}t$
    (quadratic variation, using $\sum_j x_j=M$)
  \item $d\langle x_i,M\rangle_t = \sigma_n^2 x_i\,\mathrm{d}t$
    (only the $i$-th Brownian motion is common to $dx_i$ and $dM$)
\end{itemize}
The two second-order correction terms:
\begin{align*}
  \frac{x_i}{M^3}\,d\langle M\rangle_t &= \frac{\sigma_n^2 p_i}{M}\,dt, \\
  -\frac{1}{M^2}\,d\langle x_i,M\rangle_t &= -\frac{\sigma_n^2 p_i}{M}\,dt.
\end{align*}
These \emph{cancel exactly}.  The drift of $p_i$ is therefore:
\begin{align}
  \mathrm{d}p_i
  &= -\alpha(Lp)_i\,\mathrm{d}t
   + (\text{zero-mean noise}).
  \label{eq:dp_nodrift}
\end{align}
The drift is identical to the deterministic case; the Itô correction does
\emph{not} appear in the drift of $p_i$.

\textbf{Step 2: Quadratic variation of $p_i$ and evolution of $\|p-1/n\|^2$.}
Although the drift of $p_i$ has no Itô correction (Step~1), the noise term
in \eqref{eq:dp_nodrift} has non-trivial quadratic variation.
From the noise components in Step~1, a direct calculation gives:
\begin{equation}
  d\langle p_i\rangle_t
  = \frac{\sigma_n^2 x_i(1-p_i)^2 + \sigma_n^2 p_i^2 \sum_{j\neq i} x_j}{M^2}\,dt
  = \frac{\sigma_n^2 p_i(1-p_i)}{M}\,dt,
  \label{eq:qv_pi}
\end{equation}
using $\sum_{j\neq i}x_j = M(1-p_i)$.
Let $e = p - \frac{1}{n}\mathbf{1}$.
Applying Itô's formula to $V=\|e\|^2=\sum_i(p_i-1/n)^2$:
\begin{align}
  \frac{\mathrm{d}}{\mathrm{d}t}\,\mathbb{E}[\|e\|^2]
  &= -2\alpha\,\mathbb{E}[e^\top Le]
   + \frac{\sigma_n^2}{M}\,\mathbb{E}\Bigl[\sum_i p_i(1-p_i)\Bigr].
  \label{eq:d_deviation_sq}
\end{align}
The Itô correction $\sigma_n^2\sum_i p_i(1-p_i)/M$ enters through the
\emph{quadratic variation} of $p_i$, not its drift.
Since $\sum_i p_i(1-p_i) \le 1$ and the Poincaré inequality gives
$e^\top Le\ge\lambda_2\|e\|^2$:
\[
  \frac{d}{dt}\,\mathbb{E}[\|e\|^2]
  \;\leq\; -2\alpha\lambda_2\,\mathbb{E}[\|e\|^2] + \frac{\sigma_n^2}{M_{\min}(t)}.
\]

\textbf{Step 3: Gronwall bound.}
Applying Gronwall's inequality to \eqref{eq:d_deviation_sq}:
\begin{equation}
  \mathbb{E}[\|e(t)\|^2]
  \;\le\;
  C_0\,e^{-2\alpha\lambda_2 t}
  + \frac{\sigma_n^2\,t}{M_{\min}(t)},
  \label{eq:gronwall_spectral}
\end{equation}
where $M_{\min}(t)=\min_{s\in[0,t]}\mathbb{E}[M(s)]>0$ by
Theorem~\ref{thm:wellposed} (non-negativity and $M_0>0$).
This establishes \eqref{eq:spectral_bound_corrected}.
\end{proof}

\begin{remark}[Numerical magnitude of the Ito correction]
\label{rem:ito_correction_small}
For the paper's parameter values $\sigma_n=0.015$, $T=15$, and
$M_{\min}\approx n\cdot 0.1$ (typical surviving mass), the Ito correction
at terminal time is
\begin{equation}
  \frac{\sigma_n^2\,T}{M_{\min}(T)}
  = \frac{(0.015)^2\times 15}{n\times 0.1}
  = \frac{3.375\times 10^{-3}}{n}.
\end{equation}
For $n=50$: correction $\approx 6.75\times 10^{-5}$;
for $n=100$: $3.38\times 10^{-5}$.
Meanwhile, the main spectral term at $T=15$ satisfies
$C_0\,e^{-2\alpha\lambda_2 T}\ge 0$ with $C_0\approx 0.01$--$0.1$ typically.
Thus the Ito correction is $<2\%$ of the main spectral term for all $n\ge 50$
at the paper's parameter values, confirming that the original Corollary~1
bound is quantitatively accurate: the trigger time $\tau(\lambda_2)$ computed
from the deterministic formula is a valid approximation to within $<2\%$ error
in the spectral decay.
\end{remark}

\begin{remark}[Reduction to deterministic case]
\label{rem:spectral_deterministic}
For $\sigma_n=0$ (deterministic ELAPSE), the Ito correction vanishes and
\eqref{eq:spectral_bound_corrected} reduces to the classical spectral
decay bound
$\|p(t)-\frac{1}{n}\mathbf{1}\|_2^2\le C_0\,e^{-2\alpha\lambda_2 t}$,
which is exact (not an upper bound) in the pure Laplacian case $s=0$, $\mu=0$.
The trigger time $\tau(\lambda_2)=(2\alpha\lambda_2)^{-1}\ln(C_0/\epsilon)$
is the original Corollary~1 formula, valid exactly for $\sigma_n=0$ and
approximately (within $<2\%$) for $\sigma_n=0.015$.
\end{remark}
